\section{The real place: perturbation and fixed columns}\label{sec:real}

In this section and the next we prove Theorem~\ref{thm:real}. The determinant $\Delta_N$ must be small
not only at $G$ but on an interval of length $e^{-\rho N}$ around it. Here we reduce this to determinants
in which some columns of $\cA_N$ are replaced by values of the row functions at fixed points
(\S\ref{ssec:quadrature}--\ref{ssec:andreief}). We then recall OpenAI's energy method, extended to
configurations with fewer free points (\S\ref{ssec:energy}). It splits into two cases. In the first case it
already gives the decay rate of Theorem~\ref{thm:real}. In the second case it falls short by the gain
$\Gain$ when few columns are fixed. Section~\ref{sec:signed} closes that gap. The section ends with the
trial sequences and the verification of the bounds (\S\ref{ssec:trials}--\ref{ssec:rho}).

The design is a parameter throughout. We write $\alpha_A=A/n=(b_0+q_0+2g_0)/n_0$, $\beta_h=2h_0/n_0$,
$\gamma_H=H/n=(n_0+b_0+q_0+g_0+h_0)/n_0$, $\gamma_d=d_0/n_0=(n_0-g_0)/n_0$, and
\[
 c_{07}=1-\frac{b_0+q_0+g_0}{n_0},\qquad c_{26}=3-\frac{b_0+q_0+2g_0}{n_0}.
\]
For $D$ these are $\alpha_A=11/29$, $\beta_h=3/29$, $\gamma_H=39/29$, $\gamma_d=53/58$, $c_{07}=41/58$ and
$c_{26}=76/29$.

\subsection{A quadrature formula for the coefficient of the constant}\label{ssec:quadrature}
Let $\delta=x-G$. Then $\Delta_N(x)=\det(\cA_N+\delta\cL_N)$. The question is how large the columns of
$\cL_N$ are compared with those of $\cA_N$. By \eqref{eq:L} the entries of $\cL_N$ are linear functionals of
the coefficients of the row polynomials. These polynomials have bounded degree. So the functionals can be
written exactly as finite combinations of point evaluations. If $\tau_1,\ldots,\tau_J$ are distinct points
and $J$ is at least the number of degrees of freedom of the rows, Lagrange interpolation gives
\begin{equation}\label{eq:quad}
 \cL_N(r,k)=\sum_{a=1}^JA_r(\tau_a)\,\beta_{a,k},
\end{equation}
where the weight $\beta_{a,k}$ is the value of the same functional on the $a$-th Lagrange basis polynomial.
This is a Gauss-type quadrature formula. It makes the columns of $\cL_N$ combinations of point evaluations of
the row functions that define $\cA_N$. The perturbed determinant can then be treated by the energy method.
The cost is the size of the weights.

We take $J=d_0N=53N$ nodes. Put $g_*=401/2500$ and $a_*=2g_*/(1+g_*^2)=2005000/6410801$, and let
$K=\lceil J/2\rceil$. The nodes are the points
\[
 t_{j,\pm}=\pm\sqrt{a_*^2+(1-a_*^2)(1+s_j)/2},\qquad T_K(s_j)=0,
\]
that is, the roots of $T_K\bigl((2t^2-(1+a_*^2))/(1-a_*^2)\bigr)$. We keep all $K$ positive roots and the
first $J-K$ negative ones, so that for odd $J$ exactly one root is dropped. The nodes lie in
$S=(-1,-a_*)\cup(a_*,1)$. In the coordinate $x$ with $t=2x/(1+x^2)$ every node has $|x|>g_*$. This matters
below, because the energy fields are larger near $x=0$. Dropping a root changes the nodal polynomial by a
linear factor. It changes the derivative bound and the size of the Lagrange polynomials by factors
polynomial in $N$ only. So an odd number of nodes costs nothing in the exponent. This is needed, since
$J=53N$ is odd for odd $N$.

Let $g_S$ be the Green function of $\C\setminus\bar S$ with pole at infinity. In closed form
\[
 g_S(z)=\tfrac12\operatorname{arccosh}\Bigl(\frac{|z^2-1|+|z^2-a_*^2|}{1-a_*^2}\Bigr).
\]
The Lagrange polynomials satisfy $|\ell_a(z)|\le e^{53Ng_S(z)+o(N)}$ at distance at least $c/N$ from the
nodes. This follows from the derivative identity for Chebyshev polynomials at their roots. The weight of
column $k$ is a coefficient functional of $\ell_a$ multiplied by the column filter. After removal of the
common factor of the rows it equals
\[
 \lambda_k=4\,[t^0]\,F_k(t)\,(1-t^{-2})^{-1/2},\qquad F_k(t)=t^{\,A-b-q-k}(1-t)^h(t-1)^q\ell_a(t),
\]
where $(1-t^{-2})^{-1/2}=\sum_lc_lt^{-2l}$ with $c_l=\binom{2l}l4^{-l}$. The next lemma bounds it by a
contour that depends on $k$.

\begin{lemma}[keyhole bound]\label{lem:keyhole}
Let $P$ be a real polynomial, $J\ge0$, and $F(t)=t^{-J}P(t)$. Let $0<r\le1$, and suppose that $|F|\le M$ on
the circle $|t|=r$ and $|F|\le M'$ on the segments $r\le|x|\le1$. Then
\[
 \Bigl|4\sum_{l\ge0}c_l\,[t^{J+2l}]P\Bigr|\le4(1+J)\max(M,M').
\]
\end{lemma}

\begin{proof}
Write $F=\Pi+R$ with $\Pi$ a polynomial and $R=\sum_{i=1}^J\pi_it^{-i}$ the principal part at $0$. The
functional sees only $\Pi$. Since $c_l=\frac1\pi\int_{-1}^1x^{2l}(1-x^2)^{-1/2}dx$, its value is
$\frac4\pi\int_{-1}^1\Pi(x)(1-x^2)^{-1/2}dx$, which is at most $4\max_{[-1,1]}|\Pi|$. By Cauchy's estimate
$|\pi_i|\le Mr^i$, so $|R(x)|\le JM$ for $r\le|x|\le1$, and $|\Pi|\le M'+JM$ there. For $|x|<r$ the maximum
principle gives $|\Pi(x)|\le\max_{|t|=r}|\Pi|\le(1+J)M$.
\end{proof}

In the formal proof the integral is replaced by the finite Gauss--Chebyshev identity
$\frac1M\sum_{j<M}\cos^{2l}\frac{(2j+1)\pi}{2M}=c_l$ for $l<M$, and the Cauchy estimate by the finite Cauchy
formula at roots of unity. So no integration is needed. With $\kappa=k/n$, Lemma~\ref{lem:keyhole} gives
\begin{equation}\label{eq:wk}
\begin{gathered}
 |\lambda_k|\le e^{(w(\kappa)+o(1))N},\\
 w(\kappa)=\max_{z\in\Gamma_r}\bigl[(10-58\kappa)\log|z|+6\log|1-z|+53g_S(z)\bigr],
\end{gathered}
\end{equation}
where $\Gamma_r$ is the circle $|z|=r$ together with the segments $r\le|x|\le1$. Here $10=(A-b-q)/N$ and
$6=(q+h)/N$. The radius $r$ may depend on $\kappa$. For $\kappa$ near $1$ the factor $|z|^{10-58\kappa}$
favours a large circle, and for small $\kappa$ a small one. A single circle of radius $1.001$ gives
$w(\kappa)\le52.785$ for every $\kappa$. The keyhole contour gains only a little over it at the worst
column. Its value is that the worst column for the weight is not the worst column for the binomial factor
below.

\subsection{Column expansion and an identity of Andr\'eief type}\label{ssec:andreief}
Write $\Phi(r,a)=A_r(\tau_a)$. Expanding $\det(\cA_N+\delta\Phi\beta)$ multilinearly in the columns gives a
sum over maps $f$ from the columns to $\{\varnothing,1,\ldots,J\}$ of
$\prod_{f(k)\ne\varnothing}\delta\beta_{f(k),k}\det M_f$. Here $M_f$ is the matrix whose $k$-th column is the
$k$-th column of $\cA_N$ if $f(k)=\varnothing$, and the vector $(A_r(\tau_{f(k)}))_r$ otherwise.

\begin{lemma}\label{lem:expansion}
Suppose that $|\det M_f|\le M\prod_{f(k)\ne\varnothing}C_k$ for every $f$, and that $|\beta_{a,k}|\le B_k$. Then
\[
 \bigl|\det(\cA_N+\delta\Phi\beta)-\det\cA_N\bigr|\le M\Bigl(\prod_k(1+J|\delta|B_kC_k)-1\Bigr).
\]
\end{lemma}

The lemma allows a different constant for each column. Each $M_f$ has $p=n-m$ columns that are integrals and
$m$ columns that are point values. Andr\'eief's identity extends to this situation.

\begin{lemma}\label{lem:andreief}
Let $\mu,\nu$ be finite measures, $F_r$ $(r<p+m)$, $G_i$ $(i<p)$ and $K$ integrable functions, and let $B$
be the matrix with $B(r,i)=\iint F_r(t)K(t,s)G_i(s)\,d\mu(t)\,d\nu(s)$ for $i<p$ and
$B(r,p+j)=F_r(\tau'_j)$. Then
\begin{multline*}
 (p!)^2\det B=\int\!\!\int\det\bigl[F_r((t\Vert\tau')_c)\bigr]_{r,c}\\
 \cdot\det[K(t_i,s_j)]_{i,j<p}\det[G_i(s_j)]_{j,i<p}\,d\mu^p(t)\,d\nu^p(s),
\end{multline*}
where $t\Vert\tau'$ is the list of the $p$ variables followed by the $m$ fixed points.
\end{lemma}

\begin{proof}
For $m=0$ this is Andr\'eief's identity applied twice \cite{Forrester2018Andreief}. In general, expand
$\det B$ multilinearly in its first $p$ columns. This introduces one pair of variables $(t_i,s_i)$ for each
of them. Then antisymmetrize, first in the $t$ variables and then in the $s$ variables. The last $m$ columns
are constant throughout.
\end{proof}

With $K(t,s)=1/(1-ts)$ the middle determinant is a Cauchy determinant. The third is a generalized
Vandermonde determinant $\prod_js_j^b(1-s_j)^q\det[s_j^{E_i}]$, where the exponents $E_i<n$ are those of the
columns that remain integrals. The missing exponents are the replaced columns. We bound the generalized
Vandermonde determinant by a product over the missing exponents.

\begin{lemma}\label{lem:binomvand}
Let $s_1,\ldots,s_p\in[0,1]$, let $E_1<\dots<E_p<n$ be integers, and let $\{k_1,\ldots,k_m\}$ be the
complement of $\{E_i\}$ in $\{0,\ldots,n-1\}$. Then
\[
 \bigl|\det[s_l^{E_i}]\bigr|\le\prod_{j=1}^m\beta(k_j)\prod_{i<l}|s_i-s_l|,\qquad
 \beta(k)=\sum_{i=0}^{m-1}\binom p{p-k+i},
\]
with $\binom pr=0$ for $r\notin[0,p]$.
\end{lemma}

\begin{proof}
Let $P(X)=\prod_l(X-s_l)=\sum_r(-1)^re_r(s)X^{p-r}$. Multiplying the matrix by a unitriangular matrix built
from $P$ gives $|\det[s_l^{E_i}]|=\prod_{i<l}|s_i-s_l|\cdot|\det Q|$, where $Q$ is the $m\times m$ matrix with
entries $Q_{j,i}=[X^{k_j-i}]P$. So $|Q_{j,i}|=e_{p-k_j+i}(s)\le\binom p{p-k_j+i}$, since $s\in[0,1]^p$.
Hadamard's inequality, applied to the rows of $Q$ in the $\ell^1$ norm, gives the claim.
\end{proof}

The bound does not depend on the points $s_l$. Since $\binom pr\le\binom nr$, one has
$\log\beta(k)\le n_0\tilde H(k/n)N+O(\log N)$, where $\tilde H(\kappa)$ is the binary entropy $H(\kappa)$ for
$\kappa\ge1/2$ and $\log2$ for $\kappa<1/2$. The weight \eqref{eq:wk} grows with $\kappa$ up to
$\kappa\approx0.84$, while $\tilde H$ is largest at $\kappa=1/2$. So the product of the two factors is
smaller than the product of their maxima.

\subsection{The energy estimate with fixed points}\label{ssec:energy}
OpenAI bounds the integrand in Andr\'eief's identity by the exponential of $n^2E$, where $E$ is a discrete
logarithmic energy of the points $x_i$ and $s_j$ in the coordinate $t=2x/(1+x^2)$. The energy consists of an
external field for each point, logarithmic interactions $\log|x_i-x_j|$, $\log(1-x_ix_j)$ and
$\log|s_i-s_j|$ among points of the same kind, and an interaction $-\log(1-2x_is_j+x_i^2)$ between the two
kinds. There are two cases, according to which of OpenAI's two bounds for the determinant of the row
functions is used \cite[Props.~5.2, 5.3]{OpenAI2026}. In the second case the interaction among the $x_i$ has
weight $2$ instead of $1$. In the first case the row field acquires a term with a multiplier $\lambda_1$ for
the constraint that defines the case. The second case also bounds a determinant $\det\mathsf L$ of size $n$
by $6^n$. This is the factor that Section~\ref{sec:signed} improves.

The energy is bounded from above by duality. Given trial sequences $u=(u_k)$ and $v=(v_k)$, one writes the
interaction terms as quadratic forms in the moments $\frac1n\sum_iT_k(x_i)$ and $\frac1n\sum_js_j^k$. Then one
completes the squares against $u$ and $v$. What remains is the sum of the norms $\|u\|^2=\sum_ku_k^2/k$ and
$\|v\|^2$ and, for each point, a function of one variable. With $\mathsf T(u,x)=\sum_ku_kT_k(x)/k$ and
$\mathsf S(v,x)=\sum_kv_kx^k/k$ these are, in the second case,
\begin{align*}
 X(x)&=\alpha_A\log|x|+\beta_h\log(1-x)-\gamma_H\log(1+x^2)-4\mathsf T(u,x)-\mathsf S(v,x),\\
 Y(s)&=\tfrac{b_0}{n_0}\log s+\tfrac{q_0}{n_0}\log(1-s)+2\mathsf T(v,s).
\end{align*}
In the first case they are
\begin{multline*}
 X_1(x)=\alpha_A\log|x|+\beta_h\log(1-x)-(\gamma_H-\tfrac12)\log(1+x^2)\\
 +\lambda_1\bigl(\tfrac{2g_0}{n_0}-\tfrac{2x^2}{1+x^2}\bigr)-2\mathsf T(p,x)-\mathsf S(v,x),
\end{multline*}
with the same $Y$. The trial sequence for the first variables is called $p$ in this case. If $X\le M_X$ and
$Y\le M_Y$ on the whole domains, the energy of every configuration is at most a constant determined by $M_X$,
$M_Y$ and the norms. In potential-theoretic terms the extremal configuration is described by an equilibrium
measure for two interacting measures in an external field \cite{SaffTotik}. The best trial sequences are its
moments, as Carreira observed \cite{Carreira2026}. Any trial sequences give a valid upper bound.

When $m$ of the $n$ columns are point values, the configuration has only $n-m$ free points of each kind. The
normalization is still by $n$. The fixed points enter through the determinant of the row functions evaluated
at the list $t\Vert\tau'$. They are of two kinds. Some are quadrature nodes, with $|x|>g_*$. The others are
points $y\in(0,1)$. We write $\alpha=m/n$ for the fraction of fixed columns and $\tau$ for the fraction of
fixed points of the second kind. The missing free points change the energy bound by an amount linear in
$m$. This is the energy cost of a fixed column. Every step of OpenAI's argument remains valid uniformly in
the number of free points. The regularization of the logarithmic interactions, the diagonal corrections and
the completion of squares only involve sums over the points present. For the fixed points of the first kind
the relevant field is the residual field
\[
 R(x)=-\gamma_d\log(1+x^2)-2\mathsf T(u,x),\qquad g_*\le|x|\le1 .
\]
For those of the second kind it is
\[
 N(y)=-\gamma_d\log(1+y^2)+\frac{n_0-2g_0-1}{n_0}\log y-u_1y,\qquad0<y<1,
\]
where $u_1$ is the first coefficient of $u$. The coefficient $(n_0-2g_0-1)/n_0$ is a lower bound, valid for
all $N\ge1$, of $(n-1-2g)/n$.

With Lemma~\ref{lem:binomvand} the generalized Vandermonde determinant contributes only the factor
$\prod_{i<l}|s_i-s_l|$ to the energy. The binomial factors are charged to the columns. The energy bound then
reduces to a quadratic inequality in $\alpha$ and $\tau$.

\begin{lemma}[budget, second case]\label{lem:budget2}
Let $\ell=\log2$, and let $\mathcal N,\Phi,\mathcal R,\mathcal N_c$, $c_2$ and $m\ge0$ be real numbers with
\begin{gather*}
 \mathcal N+\Phi+U-c_{07}\ell\le-m,\qquad \ell-\mathcal R+\mathcal N_c\le0,\\
 \bigl((c_{26}-1)\ell-\Phi+\mathcal R-c_2\bigr)^2\le4\ell m .
\end{gather*}
Then for all $\alpha,\tau\ge0$
\begin{multline*}
 \bigl(-(c_{07}-\tfrac12)+(c_{26}-1)\alpha-\alpha^2+\tau-\tau^2\bigr)\ell+\mathcal N+(1-\alpha)\Phi\\
 +(\alpha-\tau)\mathcal R+\tau\mathcal N_c\le-U+\tfrac12\ell+c_2\alpha .
\end{multline*}
\end{lemma}

\begin{proof}
The left side minus the right side equals
$(\mathcal N+\Phi+U-c_{07}\ell)+\alpha\bigl((c_{26}-1)\ell-\Phi+\mathcal R-c_2\bigr)-\ell\alpha^2
+\tau(\ell-\mathcal R+\mathcal N_c)-\ell\tau^2$. The last two terms are non-positive by the third hypothesis.
The first three form a concave quadratic in $\alpha$. It is non-positive for all $\alpha$, because
$a+s\alpha-\ell\alpha^2\le0$ for all $\alpha$ as soon as $s^2+4\ell a\le0$.
\end{proof}

The lemma is applied with $\mathcal N$ a bound for $2\|u\|^2+\frac12\|v\|^2$, $\Phi=M_X+M_Y$, $\mathcal R$ a
bound for the supremum of $R$ and $\mathcal N_c$ a bound for the supremum of $N$. The analogue for the first
case has no fixed points of the second kind. Its polynomial in $\alpha$ is
\[
 \bigl(-(c_{07}-\tfrac12)+(c_{26}-1)\alpha-\alpha^2\bigr)\ell+n_1+(1-\alpha)F_1+\alpha\lambda_1\tfrac{2g_0}{n_0}
 \le-U+\tfrac12\ell+c_1\alpha,
\]
with $n_1$ a bound for the norms of the first trial and $F_1=M_{X_1}+M_Y$. It holds for all real $\alpha$
when its discriminant is non-positive and $c_1\le(c_{26}-1)\ell-F_1+\lambda_1\tfrac{2g_0}{n_0}$.

\subsection{The two cases at the improved rate}\label{ssec:cases}
Theorem~\ref{thm:real} asks for the rate $U_0=U+\Gain/n_0^2$ instead of $U$. A larger target is a smaller
slack $m$ in Lemma~\ref{lem:budget2}. In the first case the trial of \S\ref{ssec:trials} has enough slack:
its budget holds at $U_0$ for every $\alpha$, with the cost $c_2$ of the second case. In the second case the
budget at $U_0$ fails for small $\alpha$ and holds for $\alpha\ge1/16$. The reason is that the constant
$U_0-U$ is absorbed by the linear term once $\alpha$ is not small. So only the second case with fewer than
$n/16$ fixed columns needs a new argument. That argument is Section~\ref{sec:signed}. It proves the bound of
the second case at $U_0$ with the same cost $c_2$, for $0\le\alpha\le1/16$.

Assume this for the moment. A determinant $M_f$ with $m$ fixed columns $k_1,\ldots,k_m$ is then at most
$C^m\prod_j\beta(k_j)e^{n_0c_2Nm}$ times $\exp\bigl(n_0^2N^2(-U_0+\frac12\log2+\varepsilon)\bigr)$, with an
absolute constant $C$. The factor $C^m$ comes from the expansion of the row functions at the fixed points.
By \eqref{eq:wk} and the bound for $\beta(k)$, the constants of Lemma~\ref{lem:expansion} satisfy
\[
 \log(B_kC_k)\le\bigl(n_0c_2+w(k/n)+n_0\tilde H(k/n)+o(1)\bigr)N\le(\rho+o(1))N,
\]
with $\rho=n_0c_2+\Psi_0$ and $\Psi_0\ge\sup_\kappa\bigl(w(\kappa)+n_0\tilde H(\kappa)\bigr)$. For
$|\delta|\le e^{-\rho N}$ Lemma~\ref{lem:expansion} gives a factor $(1+O(N^5))^n=e^{o(N^2)}$. This proves
Theorem~\ref{thm:real}, given Section~\ref{sec:signed} and the numerical bounds below.

\subsection{The trial sequences}\label{ssec:trials}
For the first case we use a trial sequence with $40$ Chebyshev coefficients and a finite number of geometric
terms, with the multiplier $\lambda_1=1.68493357$. For the second case we use trial sequences of the form
\begin{equation}\label{eq:trial}
 u_k=c_k+\operatorname{Re}\sum_jr_jz_j^k\qquad(k\ge1;\ c_k=0\text{ for }k>K),
\end{equation}
with $K=\Kcoef$ Chebyshev coefficients, which are dyadic rationals, and $\Ngeo$ geometric terms. Of these,
$\Ncomplex$ have complex $z_j$ and complex $r_j$.

The geometric terms have a simple meaning. Since $\sum_{k\ge1}z^kT_k(x)/k=-\frac12\log(1-2xz+z^2)$, a term
$(z,r)$ contributes $-\frac12\operatorname{Re}\,r\log(1-2xz+z^2)$ to the potential $\mathsf T(u,x)$. If
$x_0+i\epsilon=\frac12(z+z^{-1})$, then $1-2xz+z^2=2z(x_0+i\epsilon-x)$. So the term is, up to a constant, the
logarithmic potential of a point charge of complex strength $r$ at $x_0+i\epsilon$ and its conjugate. The
extremal measure has a gap in its support around $x=0$. Its potential has a logarithmic singularity at $x=0$
and square-root behaviour at the edges of the gap. A Chebyshev expansion converges slowly near such points.
A charge at a small distance $\epsilon$ reproduces the local behaviour with one term. We place the complex
charges at the $41$ equally spaced points $x_0=-0.16+0.008j$ $(0\le j\le40)$, all with $\epsilon=10^{-2}$.
The real terms sit at points $\pm z$ with $z$ increasing from $0.8$ to $0.999997$. The numbers $z_j$ are
rationals with at most twelve decimals. The norms of the geometric terms have closed forms. Two terms
$(z_j,r_j)$ and $(z_l,r_l)$ contribute
$-\frac12\operatorname{Re}\bigl(r_jr_l\log(1-z_jz_l)+r_j\bar r_l\log(1-z_j\bar z_l)\bigr)$ to $\|u\|^2$.

\subsection{Verification of the bounds}\label{ssec:cells}
We must bound three functions of one variable from above on their whole domains. They are the field $X$ on
$[-1,1]$, the field $Y$ on $[0,1]$, and the residual field $R$ on $g_*\le|x|\le1$. With $x=\cos\theta$, each is
a constant plus a cosine series with $K$ terms plus finitely many terms $\operatorname{Re}\,w\log(1-ae^{im\theta})$.

We cover each domain by cells $[x_c-h,x_c+h]$ whose endpoints are multiples of $2^{-40}$. On each cell we
bound the function by its Taylor expansion at $x_c$. The cells are grouped in blocks. For each block the
derivatives of the polynomial part are computed exactly at the centre of the block up to order $70$. The
remainder beyond order $70$ is bounded by a single constant for the block, multiplied by $h^{70}/70!$. The
Taylor expansion at the centre of a cell is obtained from that of the block by an exact shift. The
logarithmic terms are expanded to order $5$, with remainders bounded from the exact minimum of
$|1-ae^{im\theta}|$ on the cell. A few cells at $x=0$ and $x=\pm1$ treat the singular terms separately. The
covers of $X$, $Y$ and $R$ have \Nblocks{} blocks and \Ncells{} cells. Each cell is checked by exact rational and
integer arithmetic.

The resulting bounds are
\[
 \sup X\le\SupX,\quad \sup Y\le\SupY,\quad \sup_{g_*\le|x|\le1}R\le\ResCap .
\]
So $M_X+M_Y=\FieldCap$, and the norms satisfy $2\|u\|^2+\frac12\|v\|^2\le\NormCap$. The first coefficient
satisfies $u_1\ge9/50$, which is used near $x=0$. For the points of the second kind a one-variable estimate
shows $\sup_{0<y<1}N(y)\le-77/100$ for $u_1\ge9/50$. For the first case, with the weights
$(\alpha_A,\beta_h,\gamma_H-\frac12,\lambda_1,2g_0/n_0)=(11/29,3/29,49/58,\lambda_1,5/29)$, the bounds
$\sup X_1\le-1.15854$ and $\sup Y\le-1.32981$ are obtained by isolating the zeros of the numerators of the
derivatives. These numerators are polynomials with integer coefficients. The norm bound is $n_1=0.67761$.

With $c_{07}=41/58$, $c_{26}=76/29$ and $U=\Utarget$, the hypotheses of Lemma~\ref{lem:budget2} hold with
$m=\Mmargin$ and $c_2=\Ctwo$. The first case admits $c_1=\Cone<c_2$. At the target $U_0$ the first-case
budget still holds with the cost $c_2$, and the second-case budget holds for $\alpha\ge1/16$ with the cost
$c_2$. The trial data are fixed once and for all. The file containing them has SHA-256 digest
\begin{center}
{\footnotesize\texttt{\Digest}}.
\end{center}

\subsection{The perturbation exponent}\label{ssec:rho}
It remains to bound $\Psi_0$. For a fixed radius the function $\kappa\mapsto\max_{\Gamma_r}\Phi_\kappa$ is a
maximum of functions affine in $\kappa$, hence convex. The entropy $H$ is concave. On a cell
$[\kappa_1,\kappa_2]\subset[\frac12,1]$ we replace $H$ by its tangent at the midpoint. The sum is then
convex, and its supremum on the cell is attained at an endpoint. On $[0,\frac12]$ the term $\tilde H=\log2$
is constant. We cover $[0,1]$ by $14$ cells, each with a rational radius $r$. At each endpoint the maximum of
$\Phi_\kappa$ on the circle and on the segments is bounded by a cover of the contour by small arcs and
intervals, in exact arithmetic. This proves
\[
 \sup_{0\le\kappa\le1}\bigl(w(\kappa)+58\tilde H(\kappa)\bigr)\le\Psi_0=\PsiZero .
\]
The supremum is attained near $\kappa=0.62$, with radius about $0.6$. Hence $\rho=58\cdot\Ctwo+\PsiZero=\Rho$.
